\newpage

\subsection{Project Conclusions}

By this stage, this project constitutes the most significant challenge the author has undertaken. Prior to its initiation, certain phases addressed in this work had been previously experienced within academic coursework in collaborative team settings, as outlined in Section \ref{sec:integration-of-technical-knowledge}. However, this project marks the author's first comprehensive engagement with a complete end-to-end development process, encompassing requirements elicitation, project scheduling, budget estimation, system implementation, execution analysis, and final validation.

Throughout the project lifecycle, the author has acknowledged substantial personal and professional growth developed over four academic years at UPC. The continuous adherence to the UPC's rigorous academic standards has enabled the author to acquire the competencies required to carry out the project independently while effectively overcoming ongoing challenges.

Beyond demonstrating fundamental engineering skills and learning how software projects are managed in a professional context, it has offered a valuable opportunity to integrate and consolidate key capabilities developed throughout the author's academic career. These include self-learning, adaptability, critical thinking, problem-solving, and continuous personal development.

In conclusion, this project rises above the scope of a conventional academic assignment. It has strengthened the author's confidence in addressing challenges beyond the academic environment, fostering a greater sense of independence in a professional context and providing a clearer vision of the engineer the author aspires to become. It also reflects the level of preparedness and proficiency expected from a UPC graduate when delivering high-quality results under pressure and within constrained time and resource conditions. 
